% Options for packages loaded elsewhere
\PassOptionsToPackage{unicode=true}{hyperref}
\PassOptionsToPackage{hyphens}{url}
\PassOptionsToPackage{dvipsnames,svgnames*,x11names*}{xcolor}
%
\documentclass[11pt]{article}
\usepackage{lmodern}
\usepackage{amssymb,amsmath}
\usepackage{ifxetex,ifluatex}
\ifnum 0\ifxetex 1\fi\ifluatex 1\fi=0 % if pdftex
  \usepackage[T1]{fontenc}
  \usepackage[utf8]{inputenc}
  \usepackage{textcomp}
\else % if luatex or xelatex
  \usepackage{unicode-math}
  \defaultfontfeatures{Scale=MatchLowercase}
  \defaultfontfeatures[\rmfamily]{Ligatures=TeX,Scale=1}
\fi
\IfFileExists{upquote.sty}{\usepackage{upquote}}{}
\IfFileExists{microtype.sty}{%
  \usepackage[]{microtype}
  \UseMicrotypeSet[protrusion]{basicmath}
}{}
\makeatletter
\@ifundefined{KOMAClassName}{%
  \IfFileExists{parskip.sty}{%
    \usepackage{parskip}
  }{%
    \setlength{\parindent}{0pt}
    \setlength{\parskip}{6pt plus 2pt minus 1pt}}
}{%
  \KOMAoptions{parskip=half}}
\makeatother
\usepackage{xcolor}
\IfFileExists{xurl.sty}{\usepackage{xurl}}{}
\IfFileExists{bookmark.sty}{\usepackage{bookmark}}{\usepackage{hyperref}}
\hypersetup{
  pdftitle={AUGR: an agentic, evidence-grounded framework for review and synthesis of gene-function annotations at scale},
  colorlinks=true,
  linkcolor=blue,
  filecolor=Maroon,
  citecolor=Blue,
  urlcolor=blue,
}
\urlstyle{same}
\usepackage[margin=1in]{geometry}
\usepackage{longtable,booktabs}
\usepackage{array}
\usepackage[numbers,sort&compress]{natbib}
\IfFileExists{footnotehyper.sty}{\usepackage{footnotehyper}}{\usepackage{footnote}}
\makesavenoteenv{longtable}
\usepackage{graphicx,grffile}
\makeatletter
\def\maxwidth{\ifdim\Gin@nat@width>\linewidth\linewidth\else\Gin@nat@width\fi}
\def\maxheight{\ifdim\Gin@nat@height>\textheight\textheight\else\Gin@nat@height\fi}
\makeatother
\setkeys{Gin}{width=\maxwidth,height=\maxheight,keepaspectratio}
\setlength{\emergencystretch}{3em}
\providecommand{\tightlist}{%
  \setlength{\itemsep}{0pt}\setlength{\parskip}{0pt}}
\setcounter{secnumdepth}{4}

\newcommand{\augr}{\textsc{augr}}
\newcommand{\code}[1]{\texttt{#1}}

\title{\augr{}: an agentic, evidence-grounded framework for review and
synthesis of gene-function annotations at scale}

% NOTE: author list is a placeholder to be completed before submission.
\author{The AI4Curation Project\thanks{Correspondence: Christopher J. Mungall.
Author list is a placeholder pending finalization.}}
\date{}

\begin{document}
\maketitle

\begin{abstract}
\noindent
Functional annotation databases such as the Gene Ontology (GO) underpin almost
all high-throughput interpretation of molecular data, yet a substantial fraction
of their content is over-annotated, under-specified, or propagated electronically
without organism-specific scrutiny. Correcting this at scale is bottlenecked by
\emph{synthesis}: deciding whether a single annotation is sound requires
integrating primary literature, domain architecture, paralog and pathway context,
and orthogonal evidence --- work that expert curators do well but cannot do for
millions of annotations. We present \augr{} (Assessment via Unified Gene-evidence
Review), an agentic, evidence-grounded framework that performs this synthesis gene
by gene and records the outcome as a structured, schema-validated, human-readable
review. \augr{} couples a large-language-model agent to a typed review schema
(LinkML), a deterministic anti-hallucination validator that rejects fabricated
ontology terms and unverifiable supporting quotations, and a library of
specialized sub-pipelines for literature deep-research, ad-hoc bioinformatics,
prediction assessment, annotation-rule review, and ontology-term obsoletion. We
have applied \augr{} to a public corpus of 2{,}732 gene reviews spanning 165
organism groupings and roughly 89{,}900 existing-annotation decisions: 52\% are
accepted as core, 26\% retained as non-core, and 21\% flagged for change
(over-annotated, removed, or re-termed), with over-annotation the single largest
problem category. The same framework, applied to recent machine-learning function
predictors --- including the BioReason-Pro protein-function reasoner used here as
a running example --- shows that narrative and term-level predictions largely
restate existing domain-based pipelines and fail in systematic, biologically
diagnosable ways. \augr{} is designed as a scalable intermediate layer between
aggregate prediction benchmarks and prospective experiment, and as a production
tool for curators. All reviews, schema, validator, and pipelines are released
openly.
\end{abstract}

\section{Introduction}

Determining what each gene product does, where in the cell it acts, and in which
biological processes it participates is a central task of molecular biology and
the foundation of nearly every high-throughput analysis of omics data. The Gene
Ontology (GO) provides the shared vocabulary and the annotation corpus that make
such analyses possible \cite{ashburner2000,goconsortium2023}. GO annotations come
from two broad sources: \textbf{literature-based expert curation}, in which
professional curators read primary papers and encode findings as evidence-coded
statements; and \textbf{computational inference}, which propagates function across
the long tail of uncharacterized proteins through family signatures, phylogenetic
trees, and orthology.

Both sources accumulate error. A large literature documents systematic biases and
over-annotation in GO: experimental annotation is concentrated on a small set of
popular genes \cite{haynes2018,schnoes2013}, electronic inference is uneven in
quality \cite{skunca2012}, and the ontology and its annotations drift over time so
that legacy statements no longer reflect current understanding
\cite{gaudet2017}. The consequences are not cosmetic: enrichment analyses, pathway
reconstruction, and target identification all inherit whatever distortions are
present in the annotation base.

Correcting these problems is hard not because individual judgments are subtle in
isolation, but because each judgment requires \textbf{synthesis}. Deciding whether
``protein binding'' should be replaced by a specific adapter activity, whether a
conserved-but-degenerate enzyme domain still carries catalytic activity, or
whether a phenotype readout licenses a process annotation requires a reviewer to
pull together domain architecture, paralog subfamily context, pathway presence or
absence in the organism, genetic evidence, and orthogonal primary literature. This
is the \textbf{synthesis bottleneck}. Human experts perform it routinely, but it
does not scale to the hundreds of millions of annotations in current databases, nor
to the growing stream of machine-generated predictions awaiting triage.

The need for scalable synthesis has become especially acute with the rise of
machine-learning function predictors. Deep models and protein language models
\cite{lin2023,kulmanov2020} now emit GO terms, and a newer class of
reasoner-style systems --- exemplified by BioReason-Pro \cite{fallahpour2026},
which we use as a running example throughout this paper --- emit free-text
rationales and functional summaries alongside term lists. Aggregate benchmarks
such as CAFA \cite{radivojac2013,jiang2016,zhou2019} measure agreement with a
biased GO snapshot and operate on the bag-of-terms projection of a prediction;
they cannot say whether a model's biological rationale is mechanistically
plausible, whether it ignored organism context, or whether it merely restated an
existing domain-based inference. The most pointed recent demonstration is de
Cr\'ecy-Lagard \emph{et al.} \cite{decrecylagard2025}, who manually reviewed all
453 enzyme-function predictions made by a top-performing deep model for
uncharacterized \emph{E. coli} proteins and found only a handful genuinely novel
and correct; the rest fell into reproducible, biologically explicable error
classes --- each rejection requiring exactly the multi-evidence synthesis
described above.

We present \augr{} (\emph{Assessment via Unified Gene-evidence Review}), an
agentic curation framework that automates this synthesis at the granularity of a
single gene and records every decision as a structured, validated, human-readable
review. \augr{} is not a new predictor and not a replacement for CAFA-style
metrics; it is a complementary review layer that (i) re-examines existing curated
and electronic annotations against synthesized evidence, (ii) assesses
machine-generated predictions for biological validity rather than corpus
agreement, and (iii) feeds the resulting corrections back to ontology developers
and annotation pipelines. Its defining commitment is that an agent operating over
biological knowledge must be \emph{constrained}: every ontology term it uses and
every quotation it cites is checked against external truth, so that scale does not
come at the cost of fabrication.

This paper makes four contributions. (1) We describe the \augr{} framework: its
design principles, its typed review schema, its per-gene agentic pipeline, and the
deterministic anti-hallucination validator that makes agentic curation trustworthy
(Section~\ref{sec:framework}). (2) We describe a family of mature sub-pipelines
built on the same substrate --- existing-annotation review, prediction review,
annotation-rule review, ontology obsoletion, and ad-hoc bioinformatics ---
grounded in real curation projects (Section~\ref{sec:pipelines}). (3) We report
corpus-level results from applying \augr{} to 2{,}732 gene reviews across 165
organism groupings, quantifying the action distribution and the dominant
over-annotation patterns (Section~\ref{sec:results}). (4) We release all reviews,
the schema, the validator, and the pipelines as open resources at
\code{github.com/ai4curation/ai-gene-review}.

\section{The \augr{} framework}
\label{sec:framework}

\subsection{Design principles}

\augr{} is built around five principles that distinguish trustworthy agentic
curation from unconstrained generation.

\begin{enumerate}
\tightlist
\item \textbf{Evidence-grounded.} Every retained or proposed annotation must be
tied to a verifiable source --- a primary publication with a verbatim supporting
quotation, an explicit bioinformatic analysis, or a named upstream pipeline.
\item \textbf{Schema-constrained.} Reviews conform to a typed LinkML schema
\cite{linkml} with controlled vocabularies for actions, evidence, and assessment
categories, so that outputs are machine-validatable and aggregatable rather than
free prose.
\item \textbf{Anti-hallucination by construction.} A deterministic validator,
external to the agent, rejects fabricated ontology terms and unverifiable
quotations before a review is accepted (Section~\ref{sec:validation}).
\item \textbf{Auditable and reversible.} Each decision carries a rationale and
provenance, mirroring the transparency of family- and phylogeny-based pipelines,
so a curator can trace and overturn any call.
\item \textbf{Human-in-the-loop, not human-out.} \augr{} is designed to defer to
expert curators where evidence is incomplete (an explicit \code{UNDECIDED}
action), and its outputs are surfaced for community feedback rather than written
silently to a database.
\end{enumerate}

\subsection{The review unit and schema}

The unit of review is a single gene (and its corresponding UniProt entry),
captured as a \code{GeneReview} object. The schema's principal components are:

\begin{itemize}
\tightlist
\item a \textbf{description}: a project-independent biological summary of the gene
product;
\item \textbf{existing\_annotations}: a review of every current GO annotation
(drawn from GOA via QuickGO \cite{quickgo}), each assigned one of a small set of
controlled \emph{actions} (Table~\ref{tab:actions}) together with a reason and, where
applicable, a proposed replacement term;
\item \textbf{core\_functions}: a synthesized set of GO-CAM-like
\cite{thomas2019} ``annotons'' capturing the gene product's essential, evolved
molecular functions in context (molecular function, the process and location it
acts in, and complex membership);
\item \textbf{references}: literature and bioinformatics sources, each optionally
carrying a reviewer-supplied \code{reference\_review} that records relevance and
citation correctness (e.g.\ flagging a well-formed identifier that points to the
wrong paper);
\item \textbf{proposed\_new\_terms}, \textbf{suggested\_questions}, and
\textbf{suggested\_experiments}.
\end{itemize}

The action vocabulary is deliberately richer than accept/reject. Beyond
\code{ACCEPT} and \code{REMOVE}, \augr{} distinguishes \code{KEEP\_AS\_NON\_CORE}
(valid but peripheral to the gene's core function --- important for pleiotropic
genes), \code{MARK\_AS\_OVER\_ANNOTATED} (not wrong but over-reaching),
\code{MODIFY} (right essence, better term available --- paired with a replacement),
\code{UNDECIDED} (insufficient accessible evidence; the mandated action when the
agent cannot verify the source), and \code{NEW} (a proposed annotation not present
in GOA). Isoform-specific and negated (\code{NOT}) annotations are tracked
explicitly.

\begin{table}[htbp]
\centering
\caption{The controlled action vocabulary applied to each existing annotation.}
\label{tab:actions}
\small
\begin{tabular}{@{}>{\raggedright\arraybackslash}p{0.30\linewidth}>{\raggedright\arraybackslash}p{0.62\linewidth}@{}}
\toprule
\textbf{Action} & \textbf{Meaning} \\
\midrule
\code{ACCEPT} & Correct and well supported; retained as a core function. \\
\code{KEEP\_AS\_NON\_CORE} & Valid but not central (e.g.\ developmental or contextual roles of a pleiotropic gene). \\
\code{MARK\_AS\_OVER\_ANNOTATED} & Not strictly wrong, but an over-reach given the evidence. \\
\code{MODIFY} & Essence is sound but a better (more/less specific) term exists; a replacement is proposed. \\
\code{REMOVE} & Unlikely correct on combined evidence (e.g.\ a demonstrably wrong mapping or over-propagated inference). \\
\code{UNDECIDED} & Evidence cannot be accessed or adjudicated; defer to expert curators. \\
\code{NEW} & A function supported by evidence but absent from GOA; proposed as a new annotation. \\
\bottomrule
\end{tabular}
\end{table}

\subsection{The per-gene agentic pipeline}

For each gene, \augr{} runs a reproducible sequence of steps, most of which are
delegated to specialized agents or deterministic tools rather than to a single
monolithic prompt:

\begin{enumerate}
\tightlist
\item \textbf{Fetch.} Deterministic retrieval of the UniProt record, the GOA
annotation set, and the cited publications, plus a seeded review stub
(\code{just fetch-gene <organism> <gene>}). The corpus currently caches 20{,}690
publications and 4{,}895 Reactome entries locally.
\item \textbf{Deep research.} A literature-synthesis agent produces a cited notes
document (\code{GENE-notes.md} / \code{GENE-deep-research-*.md}) in which every
assertion is tied to a source and a verbatim supporting quotation. Multiple
research providers can be used; the resulting summaries are treated as preliminary
evidence, not ground truth.
\item \textbf{Annotation-by-annotation review.} An agent evaluates each existing
annotation against the synthesized evidence and assigns an action, a reason, and
where needed a replacement term, following GO curation guidelines (including the
True Path Rule and ECO evidence hierarchy \cite{giglio2019}). Crucially, the agent
is instructed \emph{not} to overturn an experimental annotation merely because the
cached abstract foregrounds a paralog --- the full text the curator read is often
not in cache --- defaulting instead to \code{UNDECIDED}.
\item \textbf{Ad-hoc bioinformatics (optional).} Where database lookups are
insufficient, a bioinformatics sub-agent writes and executes reproducible analysis
code (active-site residue checks, domain architecture, signal-peptide and
transmembrane prediction, cross-species comparison) in a per-gene, dependency-pinned
workspace, and reports results --- including inconclusive results --- without
fabrication.
\item \textbf{Core-function synthesis.} A synthesis agent distills the accepted and
modified annotations into a small set of GO-CAM-like core functions.
\item \textbf{Validation.} The review is checked against the schema and the
anti-hallucination validator (\code{just validate <organism> <gene>}); failures
block acceptance.
\end{enumerate}

\subsection{Anti-hallucination validation}
\label{sec:validation}

The central technical enabler of trustworthy scale is a deterministic validator
that runs \emph{outside} the agent and treats the agent's output as untrusted. It
enforces several independent checks:

\begin{itemize}
\tightlist
\item \textbf{ID/label tuple checksums.} Every ontology term must carry both an
identifier and a label, and the pair must match the canonical term in the source
ontology, resolved live via OAK/OLS \cite{oak}. A fabricated identifier, or a real
identifier with the wrong label, is rejected.
\item \textbf{Verbatim supporting text.} Every \code{supporting\_text} quotation
must be an exact substring of the cited cached publication. A paraphrased,
invented, or wrong-paper quotation fails automatically.
\item \textbf{Citation--title consistency.} Each reference's stored title must
match the fetched record, catching transposed or wrong PMIDs.
\item \textbf{Branch reachability.} Author-supplied terms in \code{core\_functions}
are bound to dynamic enums by GO aspect: a cellular-component term placed in a
molecular-function slot is a blocking error.
\end{itemize}

A deliberate design choice is \textbf{two-tier trust}: machine-sourced identifiers
that arrive from GOA (the \code{existing\_annotations} term ids) are \emph{not}
hard-validated, because they carry no hallucination risk and GOA legitimately leads
or lags the ontology release used for validation; author-supplied identifiers (the
core-function terms the agent invents) \emph{are} strictly validated. The
validator thus concentrates scrutiny exactly where fabrication can occur. Because
these checks are mechanical, they catch the citation failures that no amount of
fluent prose can hide --- and they free human review to focus on the substantive
question of whether a correctly-cited paper actually \emph{supports} the claim.

\section{Pipelines}
\label{sec:pipelines}

The same review substrate supports a family of pipelines, each matured through a
concrete curation project. We summarize the five most developed.

\subsection{Existing-annotation review}

The core pipeline re-examines the curated and electronic annotations already in
GOA. At corpus scale it surfaces recurring, source-attributable failure modes that
individual annotations hide. Several have been worked up as standalone analyses:

\begin{itemize}
\tightlist
\item \textbf{Uninformative terms.} ``Protein binding'' (GO:0005515) and similarly
generic terms are systematically flagged for \code{MODIFY} toward a specific
molecular function (e.g.\ an adapter or scaffold activity) where evidence permits.
\item \textbf{IBA over-propagation.} Phylogenetically inferred (IBA) annotations
transfer function from characterized orthologs, but break for pseudo-enzymes,
functionally diverged orthologs, and over-general parent terms --- e.g.\ histone
demethylase activity propagated to the degenerate JmjC protein Epe1, or
ubiquitin-activating activity propagated to the ISG15-specific UBA7.
\item \textbf{Enzyme specificity.} Annotations that are too narrow, too broad, or
simply the wrong reaction relative to the enzyme's demonstrated substrate range are
re-termed (e.g.\ a phosphatidylcholine-specific lysophospholipase call corrected to
the broader phospholipase B activity).
\item \textbf{Assay-to-function reliability.} A conceptual framework distinguishes
readouts that measure a gene product's own molecular activity from distal,
convergent phenotype readouts (ROS, caspase activation, the UPR, intracellular
Ca\textsuperscript{2+}), the latter being the dominant driver of process-level
over-annotation.
\item \textbf{Evidence-source sufficiency.} A meta-analysis of which cheap evidence
sources (reviews, abstracts, deep-research reports) suffice to confirm an
\code{ACCEPT}, informing how much effort each action actually requires.
\end{itemize}

\subsection{Prediction review}

A parallel schema, \code{PredictionReview}, evaluates computational predictions
that are \emph{not} in curated GOA --- from enzyme-EC predictors, InterPro2GO,
PANTHER/IBA, and reasoner-style systems. Predictions are scored not by corpus
agreement but by biological validity, using the taxonomy of de Cr\'ecy-Lagard
\emph{et al.} \cite{decrecylagard2025} (Table~\ref{tab:pred}), with explicit error
types (paralog over-annotation, ignored pathway context, in-vitro/in-vivo
ambiguity, frequency bias, training-data contamination). Applied to BioReason-Pro
\cite{fallahpour2026} over two linked benchmarks --- a 139-gene set for narrative
review and a 95-gene subset with term-level predictions --- this pipeline finds
that the model's summaries largely recapitulate InterPro2GO in narrative form,
default to ``cytosolic'' when no transmembrane domain is present (failing for
periplasmic, vacuolar, mitochondrial-matrix, and ER proteins), assume catalytic
activity from degenerate domains (pseudo-enzymes), and give near-identical generic
descriptions to distinct paralogs. As a positive control, \augr{} recapitulates
the major error calls of the de Cr\'ecy-Lagard \emph{E. coli} expert review from
an answer-key-withheld starting point \cite{mungall2026esrecolidetmini}.

\begin{table}[htbp]
\centering
\caption{Biological-validity assessment categories for predictions
(\code{PredictionReview}), after de Cr\'ecy-Lagard \emph{et al.}
\cite{decrecylagard2025}.}
\label{tab:pred}
\small
\begin{tabular}{@{}llp{0.60\linewidth}@{}}
\toprule
\textbf{Code} & \textbf{Score} & \textbf{Meaning} \\
\midrule
COR & 2 & Correct novel prediction. \\
CNN & 2 & Correct but not novel (already in training data). \\
LSP & 2 & Less precise than the existing annotation. \\
UNC & 1 & Uncertain --- cannot validate or refute. \\
PLI & 0 & Paralog incorrect (wrong subfamily). \\
NPI & 0 & Non-paralog incorrect (refuted by evidence). \\
REP & 0 & Repetition (frequency-biased prediction). \\
\bottomrule
\end{tabular}
\end{table}

\subsection{Annotation-rule review}

Automated annotation rules (UniProt's ARBA and UniRule) generate millions of GO
annotations from combinations of InterPro domains, PANTHER families, and taxonomic
constraints; a single faulty rule propagates to thousands of proteins. \augr{}
reviews rules directly, assessing parsimony, literature support, condition
overlap, GO-term specificity, and taxonomic scope, and recommends \code{ACCEPT},
\code{MODIFY}, or \code{REMOVE}. For example, a rule assigning ``adaptive
thermogenesis'' --- a mammalian physiological process --- under a broad Eukaryota
scope was caught annotating a fungal carboxypeptidase, indicating the scope should
be restricted to Mammalia. The corpus currently contains 73 rule reviews.

\subsection{Ontology obsoletion and replacement}

Review frequently reveals that the problem is the \emph{ontology}, not the
annotation: a term is mis-structured, conflates distinct concepts, or encodes a
spurious process. \augr{} tracks these as obsoletion-and-replacement campaigns
that feed corrections back to GO. The corpus includes 24 such campaigns (e.g.\
various vesicle-tethering, membrane-contact-site, and metabolic sub-pathway terms),
each documenting the offending term, the evidence against it, and a proposed
replacement or restructuring.

\subsection{Ad-hoc bioinformatics}

When neither literature nor existing annotation resolves a question, the
bioinformatics sub-agent performs targeted, reproducible analysis. Representative
cases include active-site validation that re-classifies a domain-bearing protein as
a pseudo-enzyme (Epe1's HVD-for-HXD substitution), localization analysis that
distinguishes a soluble periplasmic enzyme from a membrane-bound one, and domain
architecture comparison for small uncharacterized proteins. Analyses live in
per-gene, dependency-pinned workspaces with a \code{RESULTS.md} that can be cited
like any other reference, and the agent is explicitly prohibited from writing code
that fakes results or merely defers to a website.

\section{Results: corpus-level review}
\label{sec:results}

We applied \augr{} to a public corpus of 2{,}732 gene reviews (2{,}630 with
synthesized core functions) spanning 165 organism groupings, from human, mouse,
rat, worm, and yeast to a broad tail of bacteria, archaea, plants, fungi, and
non-model eukaryotes. Across these reviews, \augr{} adjudicated roughly 89{,}900
existing GO annotations and proposed 1{,}785 new ones.

The distribution of actions over existing annotations (Table~\ref{tab:corpus})
quantifies the over-annotation problem directly. About 52\% of existing
annotations are accepted as representing a core function and a further 26\% are
retained as valid but non-core --- so roughly four in five existing annotations are
judged biologically sound in some form. The remaining one in five is flagged for
change: \emph{over-annotation is the single largest problem category} (10\% of all
existing annotations), roughly equalling outright removal (5.5\%) and re-terming
(5.0\%) combined. Only 1\% of annotations were left \code{UNDECIDED}, indicating
that accessible evidence was usually sufficient to reach a call --- while that
escape hatch remains available precisely so the agent does not overrule curators
from incomplete evidence.

\begin{table}[htbp]
\centering
\caption{Distribution of \augr{} actions over existing GO annotations in the
public corpus ($\approx$89{,}900 decisions across 2{,}732 gene reviews; 165
organism groupings). \code{NEW} proposals (1{,}785) are tabulated separately as
they do not review an existing annotation.}
\label{tab:corpus}
\small
\begin{tabular}{@{}lrr@{}}
\toprule
\textbf{Action} & \textbf{Count} & \textbf{\% of existing} \\
\midrule
\code{ACCEPT} & 47{,}005 & 52.3 \\
\code{KEEP\_AS\_NON\_CORE} & 23{,}264 & 25.9 \\
\code{MARK\_AS\_OVER\_ANNOTATED} & 9{,}212 & 10.2 \\
\code{REMOVE} & 4{,}973 & 5.5 \\
\code{MODIFY} & 4{,}501 & 5.0 \\
\code{UNDECIDED} & 922 & 1.0 \\
\midrule
Retained (ACCEPT $+$ non-core) & 70{,}269 & 78.2 \\
Flagged for change (over-annotated $+$ remove $+$ modify) & 18{,}686 & 20.8 \\
\bottomrule
\end{tabular}
\end{table}

Two qualitative observations recur across the corpus and are consistent with the
prior literature on annotation bias \cite{schnoes2013,haynes2018}. First, the
flagged annotations are concentrated in biological-process and generic
molecular-function terms, and in a handful of catch-all terms (notably ``protein
binding''), rather than spread uniformly --- meaning that a small number of terms
and pipelines account for a disproportionate share of the problem and are
therefore high-value targets for systematic fixes. Second, the same structural
errors recur across unrelated genes: pseudo-enzyme over-calls, paralog-blind
transfer, and localization defaults. These are exactly the patterns the prediction
pipeline finds in machine-learning predictors, suggesting that current models have
learned the annotation \emph{distribution}, including its biases, rather than the
underlying biology.

The framework also yields an inverse resource: a literature-grounded register of
\textbf{function knowledge gaps} --- defensible statements of what remains unknown
about a gene product --- which complements the roughly 20\% of even well-studied
proteomes that remain functionally dark \cite{wood2019,rocha2023,stoeger2018}.

\section{Deployment and community feedback}

\augr{} reviews are not written silently into a database. Every review is rendered
as an interactive web page with per-element thumbs-up/down voting and a linked
detailed evaluation form, so that domain experts can flag disagreements at the
granularity of a single decision. The rendered corpus is browsable at
\url{https://ai4curation.io/ai-gene-review/}. This feedback is the intended channel
by which \augr{}'s automated synthesis is reconciled with expert judgment and by
which systematic errors in the framework itself are surfaced and corrected.

\section{Discussion}

\augr{} occupies a deliberately narrow niche: a scalable \emph{intermediate} layer
between aggregate function-prediction benchmarks and prospective wet-lab
experiment. Aggregate benchmarks such as CAFA \cite{radivojac2013,zhou2019} are
indispensable for tracking field-level progress but reward agreement with a biased
corpus and cannot inspect biological rationale; expert review reveals systematic
errors but does not scale. \augr{} keeps the multi-evidence synthesis that makes
expert review informative, while constraining the agent tightly enough --- typed
schema, external term and quotation validation, two-tier trust, and an explicit
defer-to-curator action --- that it can be run over tens of thousands of genes
without degenerating into fluent fabrication.

The corpus results argue that this layer is needed regardless of whether the
inputs are human- or machine-generated. One in five existing annotations is
flagged for change, with over-annotation dominating; and recent reasoner-style
predictors, despite topping aggregate benchmarks, largely restate domain-based
pipelines and fail in the same structured ways (pseudo-enzymes, paralog-blindness,
localization defaults) that \augr{} flags in the legacy corpus. The value of
\augr{} is less in any single verdict than in making these failure modes
\emph{visible, attributable, and aggregatable} --- and in feeding the corrections
back to ontology developers (obsoletion campaigns), annotation pipelines (rule
review), and model developers (prediction review).

\paragraph{Limitations.} \augr{}'s judgments are only as good as the evidence it
can access: cached publications are frequently abstract-only, and the framework is
deliberately conservative (\code{UNDECIDED}) rather than risk overruling an
experimental annotation whose full text it has not read. The action distribution
reflects the genes selected for review --- skewed toward genes where curators
suspected problems --- and should not be read as an unbiased estimate over all of
GOA. Agentic synthesis can still miss nuanced experimental detail, and the
validator guarantees that citations are \emph{real} and \emph{verbatim}, not that
they are the \emph{best} or \emph{sufficient} support for a claim --- which is why
human review and community feedback remain integral rather than optional.

\paragraph{Outlook.} The natural extensions are to widen coverage toward the dark
proteome, to close the loop from obsoletion and rule-review recommendations into
production GO releases, and to use the accumulating corpus of validated reviews as
a substrate for evaluating --- and eventually training --- the next generation of
function predictors against biological validity rather than corpus agreement.

\section*{Data and code availability}

All gene reviews, the LinkML schema, the anti-hallucination validator, the
prediction-, rule-, obsoletion-, and bioinformatics pipelines, and the project
pages are openly available at \code{github.com/ai4curation/ai-gene-review}
\cite{aigenereview}. Rendered, browsable reviews are at
\url{https://ai4curation.io/ai-gene-review/}. The \emph{E. coli} positive-control
dataset is archived on Zenodo \cite{mungall2026esrecolidetmini}.

\bibliographystyle{unsrtnat}
\bibliography{references}

\end{document}
